%% supplement.tex — Online Supplement
%% Lancet Regional Health – Americas
%% ===================================================================
\documentclass[12pt]{article}

\usepackage[utf8]{inputenc}
\usepackage[T1]{fontenc}
\usepackage{amsmath,amssymb}
\usepackage{graphicx}
\usepackage{booktabs}
\usepackage{multirow}
\usepackage{xcolor}
\usepackage{hyperref}
\usepackage[margin=2.5cm]{geometry}
\usepackage{caption}

%% Prefix tables and figures for appendix (Lancet style)
\renewcommand{\thetable}{Appendix table~\arabic{table}}
\renewcommand{\thefigure}{Appendix figure~\arabic{figure}}

\begin{document}

\begin{center}
{\Large\bfseries Online Supplement}\\[0.5em]
{\large Beyond heatwaves: the delayed respiratory burden of cold extremes
and air pollution in southern Brazil, a time-series and modelling study}\\[0.5em]
Felipe Baglioli,$^{a}$ Lucas Rover,$^{b}$ Eduardo Tadeu Bacalhau,$^{a}$
Wilson Flavio Feltrim Roseghini,$^{a}$ Camila Bertoletti Carpenedo,$^{a}$\\
Luciene Pimentel,$^{c}$ Pedro Augusto Breda Font\~{a}o,$^{a}$
Leonardo Jos\'{e} Rossoni Quadros,$^{b}$ Eloy Fassi Casagrande Jr.,$^{b}$\\
Francisco Mendon\c{c}a,$^{a}$ Yara de Souza Tadano,$^{b}$ Ricardo Henrique Moreton Godoi$^{a,*}$\\[0.3em]
{\small $^{a}$Federal University of Paran\'{a} (UFPR), Curitiba, PR, Brazil}\\
{\small $^{b}$Federal University of Technology -- Paran\'{a} (UTFPR), Ponta Grossa, PR, Brazil}\\
{\small $^{c}$Pontifical Catholic University of Paran\'{a} (PUCPR), Curitiba, PR, Brazil}
\end{center}

\bigskip

%% ═══════════════════════════════════════════════════════════════════
%% SUPPLEMENTARY METHODS
%% ═══════════════════════════════════════════════════════════════════
\section*{Supplementary Methods}

\subsection*{S.1\quad Exploratory $\text{PM}_{2{\cdot}5}$ DLNM and confounding assessment}

An exploratory DLNM for EPA-corrected $\text{PM}_{2{\cdot}5}$ ($N$=697,
lag 0--14~days, exposure df=4, lag df=4) showed no significant
cumulative associations at any percentile (appendix figure~6): P90
(10${\cdot}$2~µg/m$^3$) RR=0${\cdot}$91 [0${\cdot}$78--1${\cdot}$07].
The low ambient concentrations (median 4${\cdot}$1~µg/m$^3$, well
below the WHO 24-hour guideline) and 36${\cdot}$4\% missing data
likely limited statistical power. A joint model ($N$=683) yielded
RERI=0${\cdot}$013 [$-$0${\cdot}$006, 0${\cdot}$033] on the additive
scale, indicating no evidence of supra-additive interaction between
temperature and $\text{PM}_{2{\cdot}5}$ (appendix table~6).

\paragraph{Confounding diagnostic: A-B-C comparison (appendix table~9)}
To distinguish $\text{PM}_{2{\cdot}5}$ confounding from sample restriction
bias, we compared four models at P10 of $T_\mathrm{min}$:
\begin{itemize}
  \item \textbf{Model~A:} Unadjusted temperature DLNM on the full
    sample ($N$=1068, RR=1${\cdot}$17).
  \item \textbf{Model~B:} Unadjusted temperature DLNM restricted to
    $\text{PM}_{2{\cdot}5}$-complete days ($N$=683, RR=0${\cdot}$82) ---
    \emph{without} adjusting for $\text{PM}_{2{\cdot}5}$.
  \item \textbf{Model~C:} Temperature DLNM \emph{with} $\text{PM}_{2{\cdot}5}$
    as linear covariate ($N$=683, RR=0${\cdot}$82).
  \item \textbf{Model~D:} Temperature DLNM with relative humidity as
    linear covariate ($N$=851, RR=1${\cdot}$18).
\end{itemize}
Models~B and~C are virtually identical ($\Delta$RR$<$0${\cdot}$001),
proving that $\text{PM}_{2{\cdot}5}$ does not confound the
temperature--hospitalisation association. The entire attenuation from
Model~A to Model~B (30\% reduction) is due to restricting the sample
to $\text{PM}_{2{\cdot}5}$-available days, which systematically
underrepresent cold temperatures (appendix table~10). Model~D confirms that
humidity does not confound the primary estimates.

Atmospheric pressure was not available in the merged daily dataset and
could not be included; however, pressure's effect on respiratory
admissions is predominantly mediated through
temperature (Linares et al., 2025; ref~4 in main manuscript), and its omission is unlikely to bias
the primary temperature estimates substantially.

\subsection*{Missing data sensitivity}

\paragraph{$\text{PM}_{2{\cdot}5}$ missingness mechanism (appendix tables~10--11)}
$\text{PM}_{2{\cdot}5}$ data were available for 697 of 1096~study days
(63${\cdot}$6\%). Missingness was driven by PurpleAir sensor downtime
(hardware resets, network outages) and inter-channel quality control
exclusions. Days with missing $\text{PM}_{2{\cdot}5}$ had significantly lower
temperatures ($-$1${\cdot}$4\,°C, $p$$<$0${\cdot}$001) and higher
hospitalisations ($+$3${\cdot}$4/day, $p$$<$0${\cdot}$001; appendix table~10),
confirming that missingness is not MCAR.

To characterise the mechanism further, we fitted a logistic model:
$\mathrm{P}(\text{PM}_{2.5}\text{ observed}) \sim
\mathrm{ns}(T_\mathrm{min}, 4\,\text{df}) + \mathrm{month} +
\mathrm{DOW}$ (appendix table~11). The model achieved pseudo-$R^2$=0${\cdot}$26
and AUC=0${\cdot}$80, confirming that season and temperature
strongly predict missingness. Conditional on these covariates, no
additional pattern is apparent, consistent with a missing-at-random
(MAR) mechanism driven by seasonal sensor downtime rather than by
$\text{PM}_{2{\cdot}5}$ levels themselves.

Consequently, complete-case analyses restricted to
$\text{PM}_{2{\cdot}5}$-available days underrepresent cold temperatures and
produce attenuated temperature effect estimates. The A-B-C diagnostic
(appendix table~9) demonstrates that this attenuation is entirely attributable
to sample restriction, not $\text{PM}_{2{\cdot}5}$ confounding. The primary
temperature DLNM, which uses the full dataset ($N$=1068) without
conditioning on $\text{PM}_{2{\cdot}5}$ availability, is not affected.

\paragraph{Imputation sensitivity}
Sensitivity analyses with mean imputation and $k$-nearest-neighbour
imputation ($k$=5) produced qualitatively similar clustering solutions
to the primary complete-case analysis (silhouette scores within
0${\cdot}$02 of the main analysis). PCA loadings and explained
variance were robust to imputation method.

\subsection*{S.2\quad Clustering and PCA}

K-means clustering was applied to six standardised variables
($T_\mathrm{max}$, $T_\mathrm{med}$, $T_\mathrm{min}$,
$\text{PM}_{2{\cdot}5}$, relative humidity, atmospheric pressure) on
$N$=558 complete-case days. Optimal $K$=3 was selected by silhouette
score (0${\cdot}$357). Three clusters emerged: `Clean Cold' ($n$=248),
`Polluted Heat' ($n$=12), and `Clean Heat' ($n$=298). Kruskal--Wallis
test for between-cluster differences in daily hospitalisations
yielded $p$=0${\cdot}$039 (appendix tables~2--3). PCA on $\text{PM}_{2{\cdot}5}$,
$T_\mathrm{min}$, and hospitalisations identified PC1 (44${\cdot}$5\%
variance) associating cold, polluted days with higher admissions
(appendix table~4). Extended lag correlations (appendix table~5) revealed peak
$T_\mathrm{min}$--hospitalisation correlation at lag~9
($r$=$-$0${\cdot}$320, $p$$<$0${\cdot}$001).

\subsection*{S.3\quad Machine learning forecasting details}

\paragraph{Extreme Learning Machine (ELM)}
A single-hidden-layer feedforward network with random input weights
and regularised pseudoinverse output weights. Hyperparameter: number
of hidden neurons (scanned over 1--200).

\paragraph{Multilayer perceptron (MLP)}
A PyTorch-based feedforward network with one hidden layer (tanh
activation), trained with Adam optimiser (learning rate 0${\cdot}$01),
mini-batch size 20, early stopping (patience 10), and learning-rate
scheduling.

\paragraph{XGBoost}
A gradient-boosted tree ensemble (max depth 5, learning rate
0${\cdot}$1, \texttt{reg:squarederror} objective). The number of
estimators was tuned via validation-set MSE.

\paragraph{Training protocol}
Data were split temporally: 70\% training, 15\% validation, 15\%
test. Features were standardised (zero mean, unit variance). Model
selection used MAPE on the validation set; final evaluation used the
test set after retraining on training+validation data. SHAP (SHapley
Additive exPlanations) values were computed for the XGBoost model to
quantify feature importance.

\paragraph{Climate projections}
Future temperature scenarios were derived from the CMIP6 SSP5-8${\cdot}$5
pathway (EC-Earth3-Veg-LR model) for Curitiba's grid cell, with daily
$T_\mathrm{max}$, $T_\mathrm{med}$, and specific humidity extracted
for 2030, 2050, and 2099. Future $\text{PM}_{2{\cdot}5}$ was projected using
the most recent year's Transformer-imputed series as baseline.

\paragraph{Monte Carlo extreme-event simulations}
Heat waves were simulated by injecting random temperature perturbations
(+2--5\,°C for 3--7~days) during autumn; cold waves by analogous
decrements during spring/autumn. Wildfire $\text{PM}_{2{\cdot}5}$ spikes were
simulated by adding 70--100~µg/m$^3$ pulses (7--14~days) during
summer/spring.

%% ═══════════════════════════════════════════════════════════════════
%% SUPPLEMENTARY TABLES
%% ═══════════════════════════════════════════════════════════════════
\clearpage
\section*{Supplementary Tables}

\begin{table}[ht]
\centering
\caption{Summary of extreme temperature events, Curitiba, 2022--2024.
Heatwaves: $\geq$3 consecutive days with $T_\mathrm{max}$~$>$~TX90p;
cold spells: $\geq$3 consecutive days with $T_\mathrm{min}$~$<$~TN10p.}
\label{etab:extreme_events}
\begin{tabular}{lrrrr}
\toprule
Event type & Count & Total days & Mean duration & Range (days) \\
\midrule
Heatwaves    & 25 & 215 & 4${\cdot}$5 & 3--17 \\
Cold spells  &  4 &  39 & 4${\cdot}$0 & 3--7 \\
\bottomrule
\end{tabular}
\end{table}

\begin{table}[ht]
\centering
\caption{Cluster profiles at $K$=3 (optimal by silhouette, $N$=558
complete-case days). Temperature in °C (INMET); $\text{PM}_{2{\cdot}5}$
EPA-corrected. \textbf{Note:} The `Polluted Heat' cluster contains only
$n$=12 days ($<$3\% of the sample), limiting the statistical
reliability of its profile estimates. Results from this cluster should
be interpreted as exploratory pattern description.}
\label{etab:clusters_k3}
\begin{tabular}{lrrrrrrr}
\toprule
Cluster & $n$ & $T_\mathrm{max}$ & $T_\mathrm{med}$ &
  $T_\mathrm{min}$ & $\text{PM}_{2{\cdot}5}$ & RH (\%) & Pressure \\
  & & (°C) & (°C) & (°C) & (µg/m$^3$) & & (hPa) \\
\midrule
0 -- Clean Cold     & 248 & 21${\cdot}$3 & 16${\cdot}$2 & 12${\cdot}$7 &  4${\cdot}$8 & 83${\cdot}$4 & 914${\cdot}$6 \\
1 -- Polluted Heat  &  12 & 30${\cdot}$0 & 21${\cdot}$5 & 13${\cdot}$4 & 51${\cdot}$5 & 57${\cdot}$4 & 915${\cdot}$4 \\
2 -- Clean Heat     & 298 & 28${\cdot}$3 & 21${\cdot}$9 & 17${\cdot}$9 &  4${\cdot}$9 & 79${\cdot}$2 & 911${\cdot}$2 \\
\bottomrule
\end{tabular}
\end{table}

\begin{table}[ht]
\centering
\caption{Cluster profiles at $K$=4 (manuscript validation).
Kruskal--Wallis $H$=16${\cdot}$69, $p$$<$0${\cdot}$001 for
between-cluster differences in daily hospitalisations.}
\label{etab:clusters_k4}
\begin{tabular}{lrrrrrr}
\toprule
Cluster & $n$ & $T_\mathrm{max}$ & $T_\mathrm{med}$ &
  $\text{PM}_{2{\cdot}5}$ & RH (\%) & Hosp/day \\
  & & (°C) & (°C) & (µg/m$^3$) & & (mean $\pm$ SD) \\
\midrule
Clean Cold     & 197 & 23${\cdot}$3 & 18${\cdot}$4 &  3${\cdot}$2 & 86${\cdot}$4 & 31${\cdot}$3 $\pm$ 11${\cdot}$2 \\
Clean Heat     & 221 & 29${\cdot}$5 & 22${\cdot}$6 &  5${\cdot}$4 & 76${\cdot}$7 & 29${\cdot}$4 $\pm$ 10${\cdot}$5 \\
Polluted Heat  &  12 & 30${\cdot}$0 & 21${\cdot}$5 & 51${\cdot}$5 & 57${\cdot}$4 & 38${\cdot}$6 $\pm$ 9${\cdot}$7 \\
Moderate Trans. & 128 & 20${\cdot}$5 & 14${\cdot}$9 &  6${\cdot}$5 & 80${\cdot}$6 & 33${\cdot}$8 $\pm$ 11${\cdot}$7 \\
\bottomrule
\end{tabular}
\end{table}

\begin{table}[ht]
\centering
\caption{PCA loadings and explained variance ($N$=683 complete-case
days). Variables standardised prior to decomposition.}
\label{etab:pca}
\begin{tabular}{lcccc}
\toprule
Component & $\text{PM}_{2{\cdot}5}$ & $T_\mathrm{min}$ & Hospitalisations & Explained Var. \\
\midrule
PC1 & $+$0${\cdot}$564  & $-$0${\cdot}$580 & $+$0${\cdot}$588 & 44${\cdot}$5\% \\
PC2 & $+$0${\cdot}$810  & $+$0${\cdot}$525 & $-$0${\cdot}$259 & 28${\cdot}$2\% \\
\bottomrule
\end{tabular}
\end{table}

\begin{table}[ht]
\centering
\caption{Extended lag correlations (Pearson $r$) between exposures and
daily hospitalisations, selected lags. Bold: significant at $p$$<$0${\cdot}$05.}
\label{etab:extended_lags}
\begin{tabular}{lrrrrrr}
\toprule
 & \multicolumn{3}{c}{$\text{PM}_{2{\cdot}5}$ EPA} & \multicolumn{3}{c}{$T_\mathrm{min}$} \\
\cmidrule(lr){2-4} \cmidrule(lr){5-7}
Lag (days) & $r$ & $p$ & $N$ & $r$ & $p$ & $N$ \\
\midrule
 0 & \textbf{$+$0${\cdot}$159} & $<$0${\cdot}$001 & 697 & \textbf{$-$0${\cdot}$231} & $<$0${\cdot}$001 & 1066 \\
 3 & \textbf{$+$0${\cdot}$141} & $<$0${\cdot}$001 & 694 & \textbf{$-$0${\cdot}$306} & $<$0${\cdot}$001 & 1063 \\
 6 & \textbf{$+$0${\cdot}$165} & $<$0${\cdot}$001 & 691 & \textbf{$-$0${\cdot}$282} & $<$0${\cdot}$001 & 1060 \\
 9 &  $+$0${\cdot}$069 & 0${\cdot}$071 & 688 & \textbf{$-$0${\cdot}$320} & $<$0${\cdot}$001 & 1057 \\
13 & \textbf{$+$0${\cdot}$172} & $<$0${\cdot}$001 & 684 & \textbf{$-$0${\cdot}$232} & $<$0${\cdot}$001 & 1053 \\
21 &  $+$0${\cdot}$080 & 0${\cdot}$037 & 676 & \textbf{$-$0${\cdot}$225} & $<$0${\cdot}$001 & 1045 \\
28 &  $+$0${\cdot}$085 & 0${\cdot}$029 & 669 & \textbf{$-$0${\cdot}$183} & $<$0${\cdot}$001 & 1038 \\
\bottomrule
\end{tabular}
\end{table}

\begin{table}[ht]
\centering
\caption{Interaction analysis on the additive scale (RERI). Joint
quasi-Poisson DLNM, $N$=683.}
\label{etab:reri}
\begin{tabular}{lrrr}
\toprule
Metric & Estimate & 95\% CI lower & 95\% CI upper \\
\midrule
RERI & 0${\cdot}$013 & $-$0${\cdot}$006 & 0${\cdot}$033 \\
AP   & 0${\cdot}$019 &  ---     &  ---  \\
S    & 0${\cdot}$955 &  ---     &  ---  \\
\bottomrule
\end{tabular}
\end{table}

\begin{table}[ht]
\centering
\caption{Model diagnostics for the quasi-Poisson DLNM.}
\label{etab:diagnostics}
\begin{tabular}{lrrrrr}
\toprule
Model & $\hat{\phi}$ & DW & Ljung--Box $p$ & ACF lag\,1 & ACF lag\,7 \\
\midrule
Temperature & 1${\cdot}$322 & 2${\cdot}$004 & 0${\cdot}$137 & $-$0${\cdot}$006 & 0${\cdot}$077 \\
$\text{PM}_{2{\cdot}5}$  & 1${\cdot}$329 & 2${\cdot}$115 & 0${\cdot}$096 & $-$0${\cdot}$059 & 0${\cdot}$025 \\
\bottomrule
\end{tabular}
\end{table}

\begin{table}[ht]
\centering
\caption{Machine learning model performance for daily respiratory
hospitalisation forecasting, Curitiba, 2022--2024. Best result per
metric in bold.}
\label{etab:ml_results}
\begin{tabular}{llrrrr}
\toprule
Lag & Model & RMSE & MAE & MAPE (\%) & $N_\mathrm{param}$ \\
\midrule
0 & ELM     & 11${\cdot}$07 &  8${\cdot}$95 & 30${\cdot}$9 & 182 \\
0 & MLP     &  9${\cdot}$31 &  7${\cdot}$53 & 25${\cdot}$1 &   6 \\
0 & \textbf{XGBoost} & \textbf{8${\cdot}$37} & \textbf{6${\cdot}$30} & \textbf{19${\cdot}$4} &  29 \\
\midrule
1 & ELM     & 10${\cdot}$43 &  8${\cdot}$30 & 28${\cdot}$1 & 182 \\
1 & MLP     &  8${\cdot}$80 &  7${\cdot}$01 & 23${\cdot}$2 &   4 \\
1 & \textbf{XGBoost} & \textbf{8${\cdot}$10} & \textbf{6${\cdot}$10} & \textbf{18${\cdot}$0} &  18 \\
\midrule
5 & ELM     & 10${\cdot}$74 &  8${\cdot}$68 & 28${\cdot}$8 & 186 \\
5 & \textbf{MLP} & \textbf{8${\cdot}$14} & \textbf{6${\cdot}$09} & \textbf{19${\cdot}$1} &   3 \\
5 & XGBoost &  8${\cdot}$42 &  6${\cdot}$47 & 20${\cdot}$3 &  18 \\
\midrule
7 & ELM     & 10${\cdot}$28 &  7${\cdot}$99 & 26${\cdot}$3 &  94 \\
7 & MLP     &  9${\cdot}$67 &  7${\cdot}$60 & 24${\cdot}$6 &  12 \\
7 & \textbf{XGBoost} & \textbf{8${\cdot}$66} & \textbf{6${\cdot}$57} & \textbf{20${\cdot}$8} &  19 \\
\bottomrule
\end{tabular}
\end{table}

\begin{table}[ht]
\centering
\caption{Confounding diagnostic: A-B-C-D comparison. Cumulative RR at
P10 of $T_\mathrm{min}$ (9${\cdot}$8\,°C), centred at
MMT=17${\cdot}$8\,°C. Models B and C are virtually identical
($\Delta$RR$<$0${\cdot}$001), proving that $\text{PM}_{2{\cdot}5}$ does not
confound the temperature--hospitalisation association.}
\label{etab:adjusted_sensitivity}
\begin{tabular}{llrrrr}
\toprule
Model & Specification & RR & 95\% CI & $\hat{\phi}$ & $N$ \\
\midrule
A & Unadjusted (full sample)            & 1${\cdot}$17 & 1${\cdot}$00--1${\cdot}$37 & 1${\cdot}$32 & 1068 \\
B & Unadjusted ($\text{PM}_{2{\cdot}5}$-complete days) & 0${\cdot}$82 & 0${\cdot}$64--1${\cdot}$06 & 1${\cdot}$32 &  683 \\
C & + $\text{PM}_{2{\cdot}5}$ linear (same days as B)  & 0${\cdot}$82 & 0${\cdot}$64--1${\cdot}$06 & 1${\cdot}$32 &  683 \\
D & + Humidity linear                            & 1${\cdot}$18 & 0${\cdot}$98--1${\cdot}$42 & 1${\cdot}$27 &  851 \\
\bottomrule
\multicolumn{6}{l}{\footnotesize \textit{Key finding:} B$\approx$C ($\Delta$RR$<$0${\cdot}$001) $\Rightarrow$ $\text{PM}_{2{\cdot}5}$ is NOT a confounder.}\\
\multicolumn{6}{l}{\footnotesize A$\neq$B $\Rightarrow$ entire attenuation is from sample restriction to warmer days (appendix table~10).}
\end{tabular}
\end{table}

\begin{table}[ht]
\centering
\caption{Comparison of key variable distributions between days with
valid $\text{PM}_{2{\cdot}5}$ ($n$=697) and days with missing $\text{PM}_{2{\cdot}5}$
($n$=399). Welch $t$-test for difference in means.}
\label{etab:missing_data_comparison}
\begin{tabular}{lrrrr}
\toprule
Variable & Valid $\text{PM}_{2{\cdot}5}$ & Missing $\text{PM}_{2{\cdot}5}$ & Diff & $p$ \\
 & (mean $\pm$ SD) & (mean $\pm$ SD) & & (Welch $t$) \\
\midrule
$T_\mathrm{min}$ (°C)   & 15${\cdot}$3 $\pm$ 3${\cdot}$5 & 13${\cdot}$8 $\pm$ 3${\cdot}$9 & $+$1${\cdot}$4 & $<$0${\cdot}$001 \\
Rel.\ humidity (\%)       & 80${\cdot}$5 $\pm$ 8${\cdot}$9 & 78${\cdot}$9 $\pm$ 10${\cdot}$5 & $+$1${\cdot}$6 & 0${\cdot}$028 \\
Hospitalisations (n/day) & 31${\cdot}$4 $\pm$ 11${\cdot}$3 & 34${\cdot}$8 $\pm$ 11${\cdot}$6 & $-$3${\cdot}$4 & $<$0${\cdot}$001 \\
\bottomrule
\multicolumn{5}{l}{\footnotesize $\text{PM}_{2{\cdot}5}$ valid: $n$=697 days; missing: $n$=399 days. Missingness is not MCAR}\\
\multicolumn{5}{l}{\footnotesize (associated with colder days/higher admissions). See appendix table~11 for missingness model.}
\end{tabular}
\end{table}

\begin{table}[ht]
\centering
\caption{Logistic missingness model for $\text{PM}_{2{\cdot}5}$ data
availability. $\mathrm{P}(\text{PM}_{2.5}\text{ observed}) \sim
\mathrm{ns}(T_\mathrm{min}, 4\,\text{df}) + \mathrm{month} +
\mathrm{DOW}$, $N$=1068.}
\label{etab:missingness_model}
\begin{tabular}{lr}
\toprule
Metric & Value \\
\midrule
$N$ observed / missing & 683 / 385 \\
McFadden pseudo-$R^2$  & 0${\cdot}$256 \\
AUC (ROC)              & 0${\cdot}$802 \\
P(observed) range      & [0${\cdot}$01, 1${\cdot}$00] \\
P(observed) mean       & 0${\cdot}$64 \\
\bottomrule
\multicolumn{2}{l}{\footnotesize Season and temperature strongly predict missingness (AUC=0${\cdot}$80),}\\
\multicolumn{2}{l}{\footnotesize confirming not-MCAR. Consistent with seasonal sensor downtime.}
\end{tabular}
\end{table}

%% ═══════════════════════════════════════════════════════════════════
%% SUPPLEMENTARY FIGURES
%% ═══════════════════════════════════════════════════════════════════
\clearpage
\section*{Supplementary Figures}

\begin{figure}[ht]
\centering
\includegraphics[width=\textwidth]{../analysis/clustering/figures/fig_elbow_silhouette.pdf}
\caption{Cluster selection metrics for $K$=2--10: (a) elbow plot of
within-cluster inertia, (b) silhouette score (optimal at $K$=3),
(c) Calinski--Harabasz index, and (d) Davies--Bouldin index.}
\label{efig:k_selection}
\end{figure}

\begin{figure}[ht]
\centering
\includegraphics[width=0.85\textwidth]{../analysis/clustering/figures/fig_cluster_pca.pdf}
\caption{PCA biplot of cluster assignments ($K$=3). Loading vectors
show the contribution of each variable to PC1 (35${\cdot}$9\%) and
PC2 (22${\cdot}$4\%).}
\label{efig:cluster_pca}
\end{figure}

\begin{figure}[ht]
\centering
\includegraphics[width=\textwidth]{../analysis/clustering/figures/fig_method_comparison.pdf}
\caption{Comparison of clustering methods ($K$=3): K-means
(silhouette=0${\cdot}$357), Ward hierarchical (0${\cdot}$321), and
DBSCAN (single cluster).}
\label{efig:method_cmp}
\end{figure}

\begin{figure}[ht]
\centering
\includegraphics[width=\textwidth]{../analysis/clustering/figures/fig_lag_correlation.pdf}
\caption{Lag correlations (lags 0--14~days) between cluster membership
and daily hospitalisations: (a) Pearson $r$ and (b) Spearman $\rho$.}
\label{efig:lag_corr}
\end{figure}

\begin{figure}[ht]
\centering
\includegraphics[width=\textwidth]{../analysis/dlnm/figures/fig_dlnm_temp_3d.pdf}
\caption{(a) Contour plot and (b) 3D surface of RR across temperature
and lag dimensions.}
\label{efig:dlnm_temp_3d}
\end{figure}

\begin{figure}[ht]
\centering
\includegraphics[width=0.85\textwidth]{../analysis/dlnm/figures/fig_dlnm_pm25_overall.pdf}
\caption{Cumulative exposure--response curve for EPA-corrected
$\text{PM}_{2{\cdot}5}$ and respiratory hospitalisations (DLNM, lag
0--14~days).}
\label{efig:dlnm_pm25_overall}
\end{figure}

\begin{figure}[ht]
\centering
\includegraphics[width=0.85\textwidth]{../analysis/dlnm/figures/fig_stationarity.pdf}
\caption{Stationarity assessment: (a) observed hospitalisations with
STL trend, (b) seasonal component, (c) STL residuals (ADF
$p$=0${\cdot}$0003), and (d) lag correlations for raw vs.\ detrended
series.}
\label{efig:stationarity}
\end{figure}

\begin{figure}[ht]
\centering
\includegraphics[width=0.85\textwidth]{../analysis/forecasting_pm25/shap_summary_curitiba_lag0.png}
\caption{SHAP summary plot (beeswarm) for the XGBoost model at lag~0.
Each dot represents one day; colour indicates feature value (blue =
low, red = high). Features ranked by mean absolute SHAP value.}
\label{efig:shap}
\end{figure}

\begin{figure}[ht]
\centering
\includegraphics[width=\textwidth]{../analysis/dlnm/figures/fig_tmin_distribution_comparison.pdf}
\caption{Comparison of $T_\mathrm{min}$ distribution between the
INMET climatology (1961--2021, $n$=20\,018) and the study period
(2022--2024, $n$=1068). (a) Density histograms with P10 percentiles
marked. (b) Side-by-side boxplots. Kolmogorov--Smirnov test:
$D$=0${\cdot}$14, $p$$<$0${\cdot}$001. The study period is slightly
warmer than the full climatology (mean 14${\cdot}$8 vs 13${\cdot}$3\,°C)
but retains representation of cold extremes (P1=5${\cdot}$2 vs
1${\cdot}$4\,°C).}
\label{efig:tmin_distribution}
\end{figure}

%% ── ML forecasting: full results and climate projections ───────────

\begin{table}[ht]
\centering
\caption{Machine learning model performance across all prediction
horizons (lag 0--7~days), Curitiba, 2022--2024. Best result per lag
in bold. 70/30 train--test split.}
\label{etab:ml_full_results}
\begin{tabular}{llrrrr}
\toprule
Lag & Model & MSE & RMSE & MAE & MAPE (\%) \\
\midrule
\multirow{3}{*}{0}
  & ELM     & 127${\cdot}$1 & 11${\cdot}$27 & 9${\cdot}$14 & 31${\cdot}$5 \\
  & MLP     & 113${\cdot}$2 & 10${\cdot}$64 & 8${\cdot}$54 & 28${\cdot}$0 \\
  & \textbf{XGBoost} & \textbf{78${\cdot}$3} & \textbf{8${\cdot}$85} & \textbf{6${\cdot}$98} & \textbf{22${\cdot}$6} \\
\midrule
\multirow{3}{*}{1}
  & ELM     & 108${\cdot}$7 & 10${\cdot}$43 & 8${\cdot}$30 & 28${\cdot}$1 \\
  & MLP     &  77${\cdot}$4 &  8${\cdot}$80 & 7${\cdot}$01 & 23${\cdot}$2 \\
  & \textbf{XGBoost} & \textbf{65${\cdot}$6} & \textbf{8${\cdot}$10} & \textbf{6${\cdot}$10} & \textbf{18${\cdot}$0} \\
\midrule
\multirow{3}{*}{2}
  & ELM     & 113${\cdot}$7 & 10${\cdot}$66 & 8${\cdot}$68 & 29${\cdot}$3 \\
  & \textbf{MLP} & \textbf{79${\cdot}$3} & \textbf{8${\cdot}$90} & \textbf{7${\cdot}$08} & \textbf{22${\cdot}$3} \\
  & XGBoost &  89${\cdot}$6 &  9${\cdot}$47 & 7${\cdot}$42 & 24${\cdot}$3 \\
\midrule
\multirow{3}{*}{3}
  & ELM     & 111${\cdot}$6 & 10${\cdot}$56 & 8${\cdot}$42 & 28${\cdot}$7 \\
  & \textbf{MLP} & \textbf{69${\cdot}$5} & \textbf{8${\cdot}$34} & \textbf{6${\cdot}$31} & \textbf{20${\cdot}$0} \\
  & XGBoost &  70${\cdot}$2 &  8${\cdot}$38 & 6${\cdot}$38 & 19${\cdot}$8 \\
\midrule
\multirow{3}{*}{4}
  & ELM     & 117${\cdot}$1 & 10${\cdot}$82 & 8${\cdot}$67 & 28${\cdot}$6 \\
  & MLP     &  78${\cdot}$5 &  8${\cdot}$86 & 6${\cdot}$88 & 22${\cdot}$4 \\
  & \textbf{XGBoost} & \textbf{73${\cdot}$4} & \textbf{8${\cdot}$57} & \textbf{6${\cdot}$52} & \textbf{20${\cdot}$1} \\
\midrule
\multirow{3}{*}{5}
  & ELM     & 115${\cdot}$4 & 10${\cdot}$74 & 8${\cdot}$68 & 28${\cdot}$8 \\
  & \textbf{MLP} & \textbf{66${\cdot}$2} & \textbf{8${\cdot}$14} & \textbf{6${\cdot}$09} & \textbf{19${\cdot}$1} \\
  & XGBoost &  70${\cdot}$9 &  8${\cdot}$42 & 6${\cdot}$47 & 20${\cdot}$3 \\
\midrule
\multirow{3}{*}{6}
  & ELM     & 123${\cdot}$1 & 11${\cdot}$09 & 8${\cdot}$81 & 28${\cdot}$5 \\
  & \textbf{MLP} & \textbf{68${\cdot}$8} & \textbf{8${\cdot}$29} & \textbf{6${\cdot}$21} & \textbf{18${\cdot}$9} \\
  & XGBoost &  69${\cdot}$0 &  8${\cdot}$30 & 6${\cdot}$39 & 20${\cdot}$0 \\
\midrule
\multirow{3}{*}{7}
  & ELM     & 105${\cdot}$7 & 10${\cdot}$28 & 7${\cdot}$99 & 26${\cdot}$3 \\
  & MLP     &  93${\cdot}$5 &  9${\cdot}$67 & 7${\cdot}$60 & 24${\cdot}$6 \\
  & \textbf{XGBoost} & \textbf{75${\cdot}$0} & \textbf{8${\cdot}$66} & \textbf{6${\cdot}$57} & \textbf{20${\cdot}$8} \\
\bottomrule
\end{tabular}
\end{table}

\begin{figure}[ht]
\centering
\includegraphics[width=\textwidth]{../analysis/forecasting_pm25/Predictions_curitiba_LAG_1.png}
\caption{Observed vs.\ predicted daily respiratory hospitalisations at
the 1-day prediction horizon (lag~1, test set). From top to bottom:
ELM (MAPE=28${\cdot}$1\%), MLP (MAPE=23${\cdot}$2\%), and XGBoost
(MAPE=18${\cdot}$0\%). XGBoost most closely tracks daily variability,
including weekend troughs and post-cold-spell peaks.}
\label{efig:ml_predictions}
\end{figure}

\begin{figure}[ht]
\centering
\includegraphics[width=\textwidth]{../analysis/forecasting_pm25/comparacao_2024_2050_curitiba_lag0__.png}
\caption{Projected daily respiratory hospitalisations under CMIP6
SSP5-8${\cdot}$5 for 2050 (coloured) compared with 2024 baseline
(black), using MLP (top) and XGBoost (bottom). Baseline annual totals
are similar ($\sim$12\,300), but the 2050 scenario shows amplified
winter peaks due to projected temperature shifts.}
\label{efig:projections_baseline}
\end{figure}

\begin{figure}[ht]
\centering
\includegraphics[width=\textwidth]{../analysis/forecasting_pm25/série_quatro_ondas_2050_2050_curitiba_lag0_cold.png}
\caption{Monte Carlo cold wave simulation for 2050: daily
hospitalisations under baseline 2050 climate (black) vs.\ 2050 with
superimposed cold-spell events (coloured), for MLP (top) and XGBoost
(bottom). Cold waves produce sharp hospitalisation spikes,
particularly in autumn--winter (April--July) and late spring
(November--December).}
\label{efig:coldwave_scenarios}
\end{figure}

\begin{figure}[ht]
\centering
\includegraphics[width=\textwidth]{../analysis/forecasting_pm25/série_queimadas_2050_2050_curitiba_lag0__.png}
\caption{Monte Carlo wildfire $\text{PM}_{2{\cdot}5}$ simulation for 2050:
daily hospitalisations under baseline 2050 climate (black) vs.\ 2050
with superimposed wildfire smoke events of
70--100~µg/m$^3$ for 7--14~days (coloured), for MLP (top)
and XGBoost (bottom). Wildfire episodes generate the largest
hospitalisation surges, with weekly increases of up to 55\% (MLP) and
42\% (XGBoost).}
\label{efig:wildfire_scenarios}
\end{figure}

%% ═══════════════════════════════════════════════════════════════════
%% eTABLE 13 — Age-stratified descriptive analysis
%% ═══════════════════════════════════════════════════════════════════

\begin{table}[ht]
\centering
\caption{Age-stratified descriptive statistics and cold-sensitivity
analysis for respiratory hospitalisations, Curitiba, 2022--2024.}
\label{etab:age_stratification}
\begin{tabular}{lrrr}
\toprule
 & $<$5~years & 5--64~years & $\geq$65~years \\
\midrule
\multicolumn{4}{l}{\textit{Descriptive statistics}} \\
$N$ (admissions) & 7\,540 & 16\,186 & 11\,549 \\
Proportion (\%) & 21${\cdot}$4 & 45${\cdot}$9 & 32${\cdot}$7 \\
Daily mean (SD) & 6${\cdot}$88 (4${\cdot}$02) & 14${\cdot}$77 (6${\cdot}$95) & 10${\cdot}$54 (4${\cdot}$57) \\
In-hospital mortality (\%) & 0${\cdot}$3 & 4${\cdot}$6 & 15${\cdot}$8 \\
\midrule
\multicolumn{4}{l}{\textit{Cold-day effect ($T_\mathrm{min}$ $<$P10 vs.\ $>$P75)}} \\
Mean admissions, cold days & 8${\cdot}$84 & 16${\cdot}$17 & 11${\cdot}$84 \\
Mean admissions, warm days & 6${\cdot}$55 & 14${\cdot}$14 & 8${\cdot}$88 \\
Ratio (cold/warm) & 1${\cdot}$35 & 1${\cdot}$14 & 1${\cdot}$33 \\
\midrule
\multicolumn{4}{l}{\textit{Lagged cross-correlation ($T_\mathrm{min}$ vs.\ admissions)}} \\
Peak lag (days) & 9 & 3 & 8 \\
Peak $r$ & $-$0${\cdot}$247 & $-$0${\cdot}$189 & $-$0${\cdot}$355 \\
\midrule
\multicolumn{4}{l}{\textit{Dominant ICD-10 subgroups (\% within age group)}} \\
J20--J22 (Acute lower resp.) & 60${\cdot}$0 & 5${\cdot}$2 & 1${\cdot}$6 \\
J09--J18 (Pneumonia/Influenza) & 22${\cdot}$4 & 30${\cdot}$2 & 44${\cdot}$8 \\
J40--J47 (COPD/Asthma) & 10${\cdot}$8 & 34${\cdot}$4 & 39${\cdot}$6 \\
\bottomrule
\end{tabular}
\end{table}

%% ═══════════════════════════════════════════════════════════════════
%% eTABLE 14 — Viral confounder sensitivity analysis
%% ═══════════════════════════════════════════════════════════════════

\begin{table}[ht]
\centering
\caption{Post-hoc sensitivity analysis: effect of adjusting for
respiratory viral circulation (SIVEP-Gripe, 2022--2024) on the
$T_\mathrm{min}$--hospitalisation association, Curitiba.}
\label{etab:viral_sensitivity}
\begin{tabular}{lrr}
\toprule
Model & $R^2$ & $T_\mathrm{min}$ coefficient \\
\midrule
$T_\mathrm{min}$ only & 0${\cdot}$357 & $-$8${\cdot}$74 \\
Viral counts only (Flu + RSV + COVID) & 0${\cdot}$359 & --- \\
$T_\mathrm{min}$ + Viral counts & 0${\cdot}$599 & $-$7${\cdot}$36 \\
\midrule
\multicolumn{3}{l}{\textit{Attenuation of $T_\mathrm{min}$ effect when controlling for viruses}} \\
Coefficient attenuation & \multicolumn{2}{r}{15${\cdot}$8\%} \\
$\Delta R^2$ (adding $T_\mathrm{min}$ to viral model) & \multicolumn{2}{r}{0${\cdot}$240 ($F$=89${\cdot}$6, $p$$<$0${\cdot}$0001)} \\
\midrule
\multicolumn{3}{l}{\textit{Cold weeks ($T_\mathrm{min}$$<$10°C) vs.\ warm weeks ($T_\mathrm{min}$$>$18°C)}} \\
Mean weekly hospitalisations (cold) & \multicolumn{2}{r}{257${\cdot}$9} \\
Mean weekly hospitalisations (warm) & \multicolumn{2}{r}{193${\cdot}$0} \\
Ratio & \multicolumn{2}{r}{1${\cdot}$34 ($p$=0${\cdot}$0002)} \\
Mean viral positivity (cold) & \multicolumn{2}{r}{0${\cdot}$32} \\
Mean viral positivity (warm) & \multicolumn{2}{r}{0${\cdot}$45} \\
\midrule
\multicolumn{3}{l}{\textit{Partial correlations ($T_\mathrm{min}$ vs.\ weekly hospitalisations)}} \\
Unadjusted (lag 1~week) & \multicolumn{2}{r}{$r$=$-$0${\cdot}$598} \\
Adjusted for viral positivity & \multicolumn{2}{r}{$r$=$-$0${\cdot}$580} \\
\bottomrule
\multicolumn{3}{p{0.9\textwidth}}{\footnotesize
Data source: SIVEP-Gripe (OpenDataSUS), $n$=28\,972 SRAG
notifications in Curitiba, 2022--2024. Weekly aggregation by
epidemiological week. Viral counts: influenza (PCR or antigen
positive), RSV (PCR or antigen positive), COVID-19 (PCR, antigen,
or CLASSI\_FIN=5). $T_\mathrm{min}$ coefficient: change in weekly
hospitalisations per 1\,°C decrease. All models are OLS with weekly
data ($n$=156~weeks).}
\end{tabular}
\end{table}

%% ── NEW: Stratified DLNM table ──────────────────────────────────────
\begin{table}[ht]
\centering
\caption{Age- and diagnosis-stratified quasi-Poisson DLNMs for
$T_\mathrm{min}$ and respiratory hospitalisations, Curitiba,
2022--2024. Same model specification as the primary analysis
(lag 0--21~days, exposure df=4, lag df=5, 7~df/year seasonal spline).
Cumulative RR centred at stratum-specific MMT.}
\label{etab:stratified_dlnm}
\begin{tabular}{lrrrrrrr}
\toprule
Stratum & Mean/day & MMT (°C) & $\hat{\phi}$ &
  \multicolumn{2}{c}{P1 (5${\cdot}$2\,°C)} &
  \multicolumn{2}{c}{P10 (9${\cdot}$8\,°C)} \\
\cmidrule(lr){5-6}\cmidrule(lr){7-8}
& & & & RR & 95\% CI & RR & 95\% CI \\
\midrule
Total (all respiratory) & 32${\cdot}$1 & 17${\cdot}$7 & 1${\cdot}$40
  & 1${\cdot}$57 & 1${\cdot}$20--2${\cdot}$05
  & 1${\cdot}$19 & 1${\cdot}$01--1${\cdot}$40 \\
Children $<$5~years & 6${\cdot}$9 & 14${\cdot}$7 & 1${\cdot}$24
  & 1${\cdot}$20 & 0${\cdot}$69--2${\cdot}$10
  & 1${\cdot}$16 & 0${\cdot}$87--1${\cdot}$55 \\
Adults 5--64~years & 14${\cdot}$7 & 18${\cdot}$4 & 1${\cdot}$42
  & 1${\cdot}$75 & 1${\cdot}$16--2${\cdot}$63
  & 1${\cdot}$34 & 1${\cdot}$04--1${\cdot}$72 \\
Adults $\geq$65~years & 10${\cdot}$5 & 17${\cdot}$6 & 1${\cdot}$06
  & 1${\cdot}$63 & 1${\cdot}$10--2${\cdot}$40
  & 1${\cdot}$14 & 0${\cdot}$90--1${\cdot}$44 \\
Pneumonia (J09--J18) & 11${\cdot}$6 & 17${\cdot}$3 & 1${\cdot}$12
  & 2${\cdot}$08 & 1${\cdot}$43--3${\cdot}$01
  & 1${\cdot}$18 & 0${\cdot}$95--1${\cdot}$48 \\
COPD/asthma (J40--J47) & 5${\cdot}$7 & 15${\cdot}$6 & 1${\cdot}$14
  & 1${\cdot}$72 & 0${\cdot}$99--2${\cdot}$98
  & 1${\cdot}$26 & 0${\cdot}$93--1${\cdot}$72 \\
\midrule
\multicolumn{8}{l}{\footnotesize \textit{Heat effect (P90, 19${\cdot}$4\,°C):}} \\
Total & & & & 1${\cdot}$07 & 0${\cdot}$99--1${\cdot}$16
  & \multicolumn{2}{l}{Non-significant in all strata} \\
\bottomrule
\multicolumn{8}{p{0.95\textwidth}}{\footnotesize
$N$=1068 valid days for all strata (same $T_\mathrm{min}$ availability).
MMT=minimum morbidity temperature (stratum-specific). Bold: 95\% CI
excludes~1. Children $<$5~years and COPD/asthma lack significance
primarily due to low daily counts limiting statistical power.}
\end{tabular}
\end{table}

%% ── NEW: Sex-stratified DLNM table ─────────────────────────────────
\begin{table}[ht]
\centering
\caption{Sex-stratified quasi-Poisson DLNMs for $T_\mathrm{min}$ and
respiratory hospitalisations, Curitiba, 2022--2024. Same model
specification as the primary analysis (lag 0--21~days, exposure df=4,
lag df=5, 7~df/year seasonal spline). Cumulative RR centred at
stratum-specific MMT.}
\label{etab:sex_dlnm}
\begin{tabular}{lrrrrrrr}
\toprule
Sex & Mean/day & MMT (°C) & $\hat{\phi}$ &
  \multicolumn{2}{c}{P1 (5${\cdot}$2\,°C)} &
  \multicolumn{2}{c}{P10 (9${\cdot}$8\,°C)} \\
\cmidrule(lr){5-6}\cmidrule(lr){7-8}
& & & & RR & 95\% CI & RR & 95\% CI \\
\midrule
Male   & 16${\cdot}$1 & 11${\cdot}$6 & 1${\cdot}$24
  & 1${\cdot}$56 & 1${\cdot}$15--2${\cdot}$12
  & 1${\cdot}$04 & 0${\cdot}$97--1${\cdot}$12 \\
Female & 16${\cdot}$1 & 18${\cdot}$1 & 1${\cdot}$16
  & 1${\cdot}$69 & 1${\cdot}$20--2${\cdot}$40
  & 1${\cdot}$45 & 1${\cdot}$17--1${\cdot}$79 \\
\midrule
\multicolumn{8}{l}{\footnotesize \textit{Heat effect (P90, 19${\cdot}$4\,°C):}} \\
Male   & & & & 1${\cdot}$15 & 0${\cdot}$92--1${\cdot}$44
  & \multicolumn{2}{l}{Non-significant} \\
Female & & & & 1${\cdot}$06 & 0${\cdot}$97--1${\cdot}$16
  & \multicolumn{2}{l}{Non-significant} \\
\bottomrule
\multicolumn{8}{p{0.95\textwidth}}{\footnotesize
$N$=1068 valid days for both strata. MMT=minimum morbidity temperature
(stratum-specific). The notably lower MMT for males (11${\cdot}$6\,°C
vs.\ 18${\cdot}$1\,°C) suggests that cold effects in men manifest only
at more extreme temperatures, whereas women exhibit vulnerability
across a broader temperature range. Heat effects were non-significant
for both sexes.}
\end{tabular}
\end{table}

\end{document}
